%% sec/discussion.tex — §7 Discussion
%% Source: paper/sec7_discussion.md

\section{Discussion: Cross-Vendor Domestic-China Viability + Take-aways}
\label{sec:discussion}

\subsection{The viability question, framed cross-vendor}
\label{sec:disc-question}

Can an industrial analog-sizing flow depend on domestic-China LLM
inference? Prior work cannot answer this;
AnalogAgent~\cite{analogagent} takes the on-prem-tiny route,
leaving cloud-frontier viability open. Our four-vendor
China-frontier cohort lets us answer in three dimensions---quality,
cost, and operational maturity---and crucially to detect whether
viability is a \emph{cohort pattern} or \emph{vendor-specific}.

\subsection{Quality}
\label{sec:disc-quality}

\tbdeight\ Across the four China-frontier checkpoints, median
success rate on LC\_VCO\_20G is within $\varepsilon\%$ of the
US-frontier-3 (Opus 4.7, GPT-5.5, Gemini 2.x), with $\varepsilon =
\tbdcell$. The cross-vendor spread within China
($\tbdcell$~percentage points) is the key indicator: a narrow
spread argues the pattern is robust; a wide spread argues that
picking a specific China vendor matters more than picking
US-vs-China. Either outcome is publishable.

\subsection{Cost}
\label{sec:disc-cost}

\tbdeight\ Under the 2026-05-01 public pricing snapshot (appendix),
the China-frontier cohort lists at $\tbdcell\times$~lower
per-output-token price than the US frontier; whether this
list-price advantage survives the per-iteration reasoning-token
premium (\S\ref{sec:results-premium}) to yield a lower
\emph{per-run} cost ordering---and in particular whether DeepSeek
V4-flash and MiMo v2.5-pro fall below GPT-5.4-mini / Haiku 4.5 on
per-run USD---is the open empirical question D8 answers. The
cost-quality story for practitioners reduces to whether
$\varepsilon$-loss on quality (\S\ref{sec:disc-quality}) is
acceptable for the cost saving---answered as a per-budget operating
point in \S\ref{sec:results-ops}.

\subsection{Operational maturity}
\label{sec:disc-ops}

Three concerns appear in any LLM rollout against a real workload:
\begin{itemize}
  \item \textbf{Rate-limit / 429 behaviour} --- surfaced per LLM
    by the \S\ref{sec:failures} bucket matrix.
  \item \textbf{Reasoning-token accounting reliability} ---
    surfaced by \S\ref{sec:results-premium}.
  \item \textbf{API-surface stability} across revision pairs
    (V4-pro vs V4-flash; Sonnet 4.6 vs Haiku 4.5) --- reported
    qualitatively per vendor in the appendix.
\end{itemize}

\subsection{Take-aways}
\label{sec:disc-takeaways}

The Pareto recommendation reduces to a per-budget table
(\S\ref{sec:results-ops}): where the China-frontier cohort is
within $\varepsilon$ of US-frontier, the cost saving justifies the
choice, with operational caveats (\S\ref{sec:disc-ops}) feeding the
procurement risk decision. The reusable methodology
contribution---multi-LLM benchmark on an industrial PDK,
spanning platform (\S\ref{sec:platform}), seed-handling and
three-axis scrub audit (\S\ref{sec:method-seed},
\S\ref{sec:method-audit}), and cost-quality / failure-bucket
reporting (\S\ref{sec:results}, \S\ref{sec:failures})---outlives
the specific numbers: the $11\times$~spec-contract matrix shape
does not age.
